\section{Preliminares y notación}

Para eliminar posibles ambigüedades en el uso de ciertos términos, comenzamos presentando las variedades diferenciables y estableciendo la dirección de las cartas que será la utilizada en adelante. También conviene prestar atención a qué nos referimos cuando hablamos de sistema de coordenadas local y cómo lo denotamos.

Entendemos como \textit{variedad diferenciable n-dimensional} un par $(M,\mathcal{A})$ donde $M$ es una variedad topológica $n$-dimensional (es decir, un espacio topológico Hausdorff, segundo numerable y localmente euclídeo) y $\mathcal{A}$ un atlas de $M$. 

A su vez, un atlas de una variedad topológica $n$-dimensional $M$ será una famila $\mathcal{A}$ de pares
\begin{equation*}
    \{(\mathcal{U}_{\alpha}, ~\varphi_{\alpha}:\mathcal{U}_{\alpha}\to \mathbb{R}^n)\}_{\alpha\in\Lambda}
\end{equation*}
a los que llamaremos \textit{cartas} donde los $\mathcal{U}_{\alpha}$ forman un recubrimiento por abiertos de $M$ y las funciones $\varphi_{\alpha}$ son homeomorfismos tales que las composiciones 
\begin{equation*}
    \varphi_\beta \circ \varphi_\alpha^{-1} : \varphi_\alpha (U_\alpha \cap U_\beta)\to \varphi_\beta (U_\alpha \cap U_\beta)
\end{equation*}
son de clase $C^\infty$ para cualesquiera $\alpha$ y $\beta$.

Todo espacio $M$ denotará una variedad diferenciable $n$-dimensional salvo que se especifique lo contrario.

\bigskip

Con frecuencia hablaremos de un \textit{sistema de coordenadas local en un entorno $\mathcal{U}$ de $p\in M$}. Con esto no nos referiremos a otra cosa que a la elección de una carta. Lo denotaremos $(x_{1}, \dots ,x_{n})$ y quedará implícito que 
\begin{equation*}
    x_i = \pi_i \circ \varphi_{\alpha} : \mathcal{U}\cap\mathcal{U}_{\alpha} \to \mathbb{R},
\end{equation*}
donde $\alpha$ dependerá del punto de $\mathcal{U}$ sobre el que actúe $x_i$ y la carta que lo contenga escogida. En términos más sencillos, $x_i$ toma un punto y lo lleva a la proyección sobre la $i$-ésima coordenada de su imagen por una de las cartas.

Haremos abuso de la notación para refererirnos indistintamente como $f$ a una función $f:M\to \mathbb{R}$ y su expresión en coordenadas $\tilde{f} = f\circ\varphi_{\alpha}^{-1} : \mathbb{R}^n \to \mathbb{R}$. Es decir, hablaremos de la imagen de un punto $p$ por $f$ como $f(p)$ y como $f(x_1,\dots,x_n)$, donde $(x_1,\dots,x_n)=(x_1(p),\dots,x_n(p))$ son sus coordenadas en una carta escogida. 

\bigskip

Por otro lado, una \textit{aplicación diferenciable} $f \colon M \to N$ entre dos variedades diferenciables es una aplicación continua tal que, para todo par de cartas
$(U,\varphi \colon U \to V) \in \mathcal A$ y $(U',\varphi' \colon U' \to V') \in \mathcal A'$, la aplicación
\[
\varphi' \circ f \circ \varphi^{-1}
\;:\;
\varphi\bigl(f^{-1}(U')\bigr) \subseteq V
\longrightarrow V'
\]
es de clase $C^\infty$. En lo que sigue, nos referiremos a ellas como aplicaciones \textit{suaves} y \textit{difeomorfismos} si poseen inversa diferenciable.

\bigskip

Cuando tengamos una aplicación diferenciable $f$ podremos hablar de su \textit{diferencial en un punto} $p \in M$, que será la aplicación lineal
\[
\mathrm d_p f \colon T_p M \longrightarrow T_{f(p)} N
\]
definida por
\[
\mathrm d_p f(X) = X(\,\bullet\, \circ f),
\]
donde $\bullet$ indica la posición del argumento. De manera explícita, el vector $\mathrm d_p f(X) \in T_{f(p)}N$ es la derivación  que, para toda función $g$ diferenciable en un entorno de $f(p)$, viene dada por
\[
\mathrm d_p f(X)(g) = X(g \circ f).
\]

Observemos que esta diferencial será la aplicación lineal habitual dada por la matriz de las derivadas parciales de $f$ cuando $f$ sea una función de $\mathbb{R}^n$ en $\mathbb{R}^m$.

\bigskip

Dada una variedad diferenciable $M$ llamaremos \textit{fibrado tangente} de $M$ a la unión disjunta
de los espacios tangentes en cada uno de los puntos de $M$, es decir,
\[
TM \;=\; \bigsqcup_{p \in M} T_pM.
\]

A su vez, un \textit{campo vectorial suave} sobre una $M$ es una aplicación suave
\[
X \colon M \to TM,
\]
que denotamos $p \mapsto X_p$, y que verifica que, para todo $p \in M$, se tiene $X_p \in T_pM$.

\bigskip

Por último, hablaremos de \textit{variedades Riemannianas} cuando lo que tengamos sea una variedad diferenciable $M$ equipada con un \textit{producto escalar}, es decir, un par $(M,\langle\cdot,\cdot\rangle)$ donde $M$ es una variedad diferenciable y para cada punto $p \in M$, $\langle\cdot,\cdot\rangle_p$ es un producto escalar (forma bilineal, simétrica y definida positiva)
  \[
    \langle\cdot,\cdot\rangle_p \colon T_pM \times T_pM \to \mathbb{R},
  \]
que verifica que para cualesquiera campos vectoriales suaves $X$ e $Y$, la función
  \[
    p \mapsto \langle X(p), Y(p) \rangle_p
  \]
es suave sobre $M$. Este producto esacalar se conoce como \textit{métrica de Riemann}. Cabe mencionar que en ocasiones nos referiremos a una variedad riemanniana solamente como $M$ si no necesitamos hacer uso explícito de la métrica ya que toda variedad diferenciable se puede equipar con una.

\bigskip

Para mayor profundidad en la mayoría de estos conceptos sugerimos que se recurra a \cite{GonzalezPrietoVariedades}, de donde nosotros mismos hemos obtenido gran parte de la información. 

Ahora que nos hemos deshecho de posibles lagunas en lo que a las variedades diferenciables se refiere abrámonos paso a la \textit{teoría de Morse}.